\section{Discussion}
There are a few things to discus. 

If you increase the size of the lattice, more skyrmions can be formed, meaning that the statistics can be done better. Because now a new lattice is generated 4 times for the same combination of parameters and checked but because we are not even close to the thermodynamical limit the statistics from these runs deviate a lot giving the large error bars in figures~\ref{} and \ref{}.


The detection program doesn't always detect the same number of skyrmions as the wrapping number suggest meaning that the detection algorithm can be improved.


Secondly the method between distinguishing from a skyrmion and an thermal fluctuation is not perfect, we see results comparable to other papers however the hight of the negative peak being 1.5 times as big is an arbitrary choice.


We expect very few positive skyrmions to be formed compared to negative skyrmions for a positive H, however this has not been tested. If this is not the case then the spiral domain may actually be a skyrmion phase where there are many positive and negative skyrmions.

As Sown in appendix~\ref{app:Temp graphs}, with temperature the skyrmionic phase in the formation analysis is gets smaller with higher temperatures, this is however not what we see in figure~\ref{fig: Formation whole domain}. We see that the formation temperature for the domain where skyrmions persist in our criteria has the lowest formation temperature. Meaning that somewhere in the code for detecting the skyrmionic phase there is an error. Due to time restraints this error has not been found and is recommended to be looked at in further studies.

The destruction temperature behaves as would be expected with this method.


\textbf{can do exponentail decay fit on the data to get the constant instead of doing the 30\% rule}